% ----------------------------------------------------
% Introduction
% ----------------------------------------------------
\documentclass[class=report,11pt,crop=false]{standalone}
\input{../Style/ChapterStyle.tex}
\makenoidxglossaries

\newacronym{fm}{FM}{frequency modulation}
\newacronym{lora}{LoRa}{Long Range}
\newacronym{iot}{IOT}{Internet of Things}
\newacronym{fsk}{FSK}{Frequency Shift Keying}
\newacronym{chirp}{CHIRP}{Compressed High Intensity Radar Pulse}
\newacronym{atps}{ATPs}{Acceptance Test Protocols}
\newacronym{hal}{HAL}{Hardware Abstraction Layer}
\newacronym{oled}{OLED}{Organic Light Emitting Diode}
\newacronym{rssi}{RSSI}{Received Signal Strength Index}
\newacronym{uart}{UART}{Universal Asynchronous Receive and Transmit}
\newacronym{usb}{USB}{Universal Serial Bus}
\newacronym{css}{CSS}{Chirp Spread Spectrum}
\newacronym{i2c}{I2C}{Inter-Integrated Circuit}
\newacronym{spi}{SPI}{Serial Peripheral Interface}
\newacronym{gps}{GPS}{Global Positioning System}
\newacronym{nmea}{NMEA}{National Marine Electronics Association}
\newacronym{mac}{MAC}{Media Access Control}
\newacronym{sf}{SF}{Spreading Factor}
\newacronym{arr}{ARR}{Auto Reload Register}
\newacronym{sdr}{SDR}{Software Defined Radio}
\newacronym{pcb}{PCB}{Printed Circuit Board}
\newacronym{fdm}{FDM}{Frequency Division Multiplexing}


\newacronym{radar}{RADAR}{Radio Detection and Ranging}
\newacronym{ofdm}{OFDM}{Orthogonal Frequency Division Multiplexing}
\newacronym{tdm}{TDM}{Time Division Multiplexing}
\newacronym{tdma}{TDMA}{Time Division Multiple Access}
\newacronym{att}{AT\&T}{American Telephone and Telegraph Company}
\newacronym{usrp}{USRP}{Universal Software Radio Peripheral}
\newacronym{uct}{UCT}{University of Cape Town}


\begin{document}
% ----------------------------------------------------
\chapter{Introduction}\label{ch:Introduction}
%\epigraph{Everyday we wake up in a world of infinite possibility. %Therefore boredom cannot exist.}%
%    {\emph{---Daniel Jones}}
%\vspace{0.5cm}
% ----------------------------------------------------
This report outlines a novel use of the \gls{ofdm} modulation technique as the backbone of the waveform for a \gls{radar} 
system as an aid to object classification. Some research into OFDM and RADAR as a whole is presented, followed by the design 
of the transmission system and the waveform itself, and then validation of the system is outlined, finally, tests are performed 
to determine the realistic performance of the system as a whole at classifying objects. In this introduction chapter, some 
background information is presented, followed by a problem statement, objectives, system requirements, and scope and limitations.

\section{Background}
RADAR technology has unlocked the ability for humans to be able to see the world around us with a new perspective. 
This ability to be able to see objects and more specifically, to see information such as position and velocity about 
these objects has helped with the development of devices and services that are of paramount importance to society. 
All the current use cases of RADAR do however lack the ability to determine exactly what object is being seen by the 
overall RADAR system. That is where the use of the OFDM waveform in a RADAR system becomes useful.



\section{Problem Statement}
Current RADAR technology works very well when determining the distance and direction of an object, however, it does 
not cope well with specific object classification tasks. One of the reasons for this is that current RADAR utilize 
either a pulse, continuous wave, or chirp waveform, which does not allow for object detection through statistical 
analysis due to the simple nature of the waveform. This is a drawback of using RADAR systems, since if object 
classification is desired, then a combination of radar and some other technology such as cameras would have to 
be used which increases the cost and complexity of the system as a whole.

\section{Objectives}
The goal of this project is to build a RADAR system that can simultaneously act as a normal RADAR which can determine 
the distance and direction of an object, as well as classify the specific object that it is looking at when compared 
to other example objects. The aim is to implement an OFDM based waveform into a RADAR system that would give the 
flexibility to analyze the returned signal statistically and determine what exactly the object is. The goal is 
to make a non-realtime system, that can distinguish between three different objects.


\section{System Requirements}
Since this system is very experimental in nature, the requirements are quite broad as of writing this, at the 
beginning of the project \colorbox{BurntOrange}{OBVIOUSLY NEED TO CHANGE THIS EVENTUALLY}. The system needs to 
perform the following tasks.

\begin{enumerate}
    \item Be able to operate as a normal RADAR. This entails determining range, azimuth, and elevation. 
    \item Be able to classify between three separate reflectors/scatterers.
\end{enumerate}

\section{Scope and Limitations}
The scope of this project has been set to not include design of RADAR hardware specifically. This is because 
design of RADAR hardware is a complex task that would probably encompass an entire masters project in and of 
itself. Therefore, the RADAR hardware portion of the project will be handled by a \gls{sdr}. This extrapolates 
the complec RADAR hardware design aspect away from the project and allows me to focus on the data collection and 
processing aspect of the project.

% ----------------------------------------------------
\ifstandalone
\bibliography{Bibliography/References}
\printnoidxglossary[type=\acronymtype,nonumberlist]
\fi
\end{document}
% ----------------------------------------------------